\chapter{系统灵敏度与稳定性研究}

\section{研究目的}
灵敏度研究用于回答两个问题：第一，观测噪声变化会如何影响位置估计；第二，测站布局变化会如何影响解的稳定性。对于工程系统，这两类问题直接关系到设备选型与站位规划，因此有必要在课程模型层面进行量化分析。

\section{噪声灵敏度指标}
设基准误差方差为 $(\sigma_A,\sigma_\rho)$，在放大系数 $\alpha$ 下变为 $(\alpha\sigma_A,\alpha\sigma_\rho)$。定义相对灵敏度指标：
\[
S_x(\alpha)=\frac{RMSE_x(\alpha)-RMSE_x(1)}{RMSE_x(1)},
\]
$S_y(\alpha),S_z(\alpha)$ 同理。若 $S_z(\alpha)$ 增长明显快于 $S_x,S_y$，说明系统垂向鲁棒性较弱。

\section{站位几何灵敏度指标}
设站位矩阵为 $\mathcal Q=\{\bm q_i\}$，定义几何扰动后位置误差增量
\[
\Delta E(\delta\mathcal Q)=
\|\hat{\bm p}(\mathcal Q+\delta\mathcal Q)-\hat{\bm p}(\mathcal Q)\|_2.
\]
进一步统计均值与分位数，可以评估站位误差传播能力。工程上可将该指标用于站位容差设计。

\section{实验分组设计}
本文设置四类实验组：
\begin{enumerate}
  \item \textbf{组 A}：仅放大测角噪声，固定测距噪声；
  \item \textbf{组 B}：仅放大测距噪声，固定测角噪声；
  \item \textbf{组 C}：测角测距同时放大；
  \item \textbf{组 D}：固定噪声，扰动测站坐标。
\end{enumerate}
每组进行 300 次蒙特卡洛重复，统计均值、标准差、95\% 分位数。

\section{结果讨论：测角噪声影响}
当仅增大测角噪声时，实验现象如下：
\begin{itemize}
  \item 水平分量误差增幅更明显；
  \item 法方程法与牛顿法误差趋势一致；
  \item 初值较优时两种方法失败率均接近 0。
\end{itemize}
这说明角度信息主要决定方向分辨率，在远距条件下会显著影响平面定位精度。

\section{结果讨论：测距噪声影响}
当仅增大测距噪声时：
\begin{itemize}
  \item 三轴误差均增大，但垂向误差增长率最高；
  \item 法化矩阵第三主轴方差增幅明显；
  \item 置信椭球长轴快速拉长，反映高度方向可观性下降。
\end{itemize}
该结果与前述 DOP 指标分析一致。

\section{结果讨论：站位扰动影响}
在站位扰动试验中，若测站整体保持环形对称结构，系统对小扰动相对不敏感；若扰动导致测站局部聚集，则条件数显著增大，收敛速度下降。此结论说明站位“均匀分布”优于“局部集中”。

\section{失败样本特征}
对极少数失败样本进行回溯，常见原因包括：
\begin{enumerate}
  \item 初值落在角度函数导数突变区附近；
  \item 法化矩阵接近奇异，步长过大；
  \item 方位角跨界未正确修正导致残差异常。
\end{enumerate}
引入阻尼项和线搜索可有效降低失败率。

\section{稳定性提升策略}
基于灵敏度结论，提出以下策略：
\begin{itemize}
  \item 使用多初值并行求解，按目标函数值选优；
  \item 对角度残差进行连续化处理；
  \item 在法方程中加入自适应阻尼；
  \item 对几何差的测站施加惩罚权重或剔除。
\end{itemize}

\section{统计表汇总}
\begin{table}[H]
\centering
\caption{噪声放大下 RMSE 变化（示意）}
\begin{tabular}{cccccc}
\toprule
$\alpha$ & 方法 & $RMSE_x$ & $RMSE_y$ & $RMSE_z$ & 失败率 \\
\midrule
0.5 & 牛顿法 & 0.010 & 0.011 & 0.121 & 0.0\% \\
1.0 & 牛顿法 & 0.015 & 0.016 & 0.175 & 0.0\% \\
2.0 & 牛顿法 & 0.028 & 0.030 & 0.335 & 0.3\% \\
4.0 & 牛顿法 & 0.056 & 0.060 & 0.672 & 1.7\% \\
\midrule
0.5 & 法方程法 & 0.011 & 0.012 & 0.126 & 0.0\% \\
1.0 & 法方程法 & 0.016 & 0.017 & 0.181 & 0.0\% \\
2.0 & 法方程法 & 0.030 & 0.032 & 0.347 & 0.5\% \\
4.0 & 法方程法 & 0.061 & 0.064 & 0.701 & 2.1\% \\
\bottomrule
\end{tabular}
\end{table}

\section{本章小结}
灵敏度研究表明：噪声增长和几何劣化都会显著损害定位性能，其中垂向分量更容易受影响。合理站位与稳健求解策略是提升系统实用性的关键。
